% Source: https://artofproblemsolving.com/community/c3h2043104_an_easy_rectangular_prism_question
% Stable identity: aops-topic-2043104-q01
Let the space diagonal be \(PQ\). Go from \(P\) along the diagonal of the \(24\times8\) base first, then up the edge of length \(6\).

\begin{asy}
usepackage("amsmath");
size(6cm);
defaultpen(fontsize(10pt));
pair ex=(1,0), ey=(0.45,0.3), ez=(0,1);
pair P(real x,real y,real z){return (x/4)*ex+(y/4)*ey+(z/4)*ez;}
real a=24,b=8,c=6;
draw(P(0,0,0)--P(a,0,0)--P(a,0,c)--P(0,0,c)--cycle,linewidth(0.9));
draw(P(a,0,0)--P(a,b,0)--P(a,b,c)--P(a,0,c)); draw(P(0,0,c)--P(0,b,c)--P(a,b,c));
draw(P(0,0,0)--P(0,b,0)--P(a,b,0),dashed); draw(P(0,b,0)--P(0,b,c),dashed);
draw(P(0,0,0)--P(a,b,0),dashed+gray); draw(P(0,0,0)--P(a,b,c),red+linewidth(1));
label("$24$",P(a/2,0,0),S); label("$8$",P(a,b/2,0),SE); label("$6$",P(0,0,c/2),W);
label("$P$",P(0,0,0),SW); label("$Q$",P(a,b,c),NE);
label("\mbox{}",(min(currentpicture,user=true)+max(currentpicture,user=true))/2,black); // keeps the text after the figure black when LaTeX typesets the labels inline
\end{asy}

The base diagonal has length \(\sqrt{24^2+8^2}=\sqrt{640}\), and it is perpendicular to the vertical edge of length \(6\), so by the Pythagorean theorem
\[
PQ=\sqrt{24^2+8^2+6^2}=\sqrt{576+64+36}=\sqrt{676}=\boxed{26}.
\]
